\section{Introduction}\label{sec:introduction}

A clique minor records a form of complete interaction inside a graph:
its vertices are represented by disjoint connected sets, and every two
sets have an edge between them. For a finite nonempty simple graph $G$,
let $h(G)$ be the largest order of a clique minor in $G$. Hadwiger
conjectured that the chromatic number always satisfies
$\chi(G)\le h(G)$ \cite{Hadwiger1943}. The companion manuscript
\cite[Corollary~1.2]{OpenAIHadwigerCounterexample2026} constructs finite
simple graphs with $\chi(G)>h(G)$, disproving this ordinary coefficient-one
assertion. Its linear relaxation asks whether
an absolute constant can replace the coefficient one. We prove a bound
of this form even when allowable colors vary from vertex to vertex.

A \emph{list assignment} $\mathcal L$ assigns a finite set of colors to
each vertex. An $\mathcal L$-coloring is a proper coloring $c$ satisfying
$c(v)\in\mathcal L(v)$ at every vertex. The list chromatic number
$\ell(G)=\chi_{\mathrm{list}}(G)$ is the least integer $k$ for which every
assignment with $|\mathcal L(v)|\ge k$ is colorable. List coloring was
introduced independently by Vizing \cite{Vizing1976} and by Erd\H{o}s,
Rubin, and Taylor \cite{ErdosRubinTaylor1980}. The linear list
version of Hadwiger's conjecture was formulated by Kawarabayashi and
Mohar \cite[Conjecture~1.3]{KawarabayashiMohar2007}.

\begin{theorem}\label{thm:main}
There is an absolute integer $C\ge1$ such that every finite nonempty
simple graph $G$ satisfies
\[
 \ell(G)\le C h(G).
\]
\end{theorem}

Thus the Linear List Hadwiger conjecture has an affirmative resolution.
The constant is independent of the graph order, the list assignment, and
the number of distinct colors appearing in its union. Taking the same
list at every vertex immediately gives the ordinary conclusion.

\begin{corollary}\label{cor:ordinary}
Every finite nonempty simple graph $G$ satisfies $\chi(G)\le C h(G)$,
with the constant $C$ in Theorem~\ref{thm:main}.
\end{corollary}
\begin{proof}
A coloring from the common list $\{1,\ldots,\ell(G)\}$ shows that
$\chi(G)\le\ell(G)$.
\end{proof}

\subsection{Prior work and the two obstacles}

The density theorems of Kostochka \cite{Kostochka1982,Kostochka1984}
and Thomason \cite{Thomason1984} show that every $K_t$-minor-free graph
has a vertex of degree $O(t\sqrt{\log t})$, as does each of its
subgraphs. Greedy coloring in reverse deletion order therefore gives
this bound for arbitrary lists. Norin, Postle, and Song
\cite{NorinPostleSong2023} went beyond that degeneracy threshold for
ordinary coloring, proving $O(t(\log t)^\beta)$ for every fixed
$\beta>1/4$. Norin and Postle \cite{NorinPostle2023} extended such
bounds to list coloring using the structure of highly connected
minor-free graphs.

The coefficient in a linear list bound has a different role from the
coefficient in ordinary Hadwiger's conjecture. Voigt
\cite{Voigt1993} constructed planar graphs that are not $4$-choosable.
Bar\'at, Joret, and Wood \cite{BaratJoretWood2011} constructed graphs
with no $K_{3s+2}$ minor that are not $4s$-choosable for every integer
$s\ge1$, and Steiner
\cite{Steiner2022List} proved that for every $\varepsilon>0$ and all
sufficiently large $t$ there are $K_t$-minor-free graphs with list
chromatic number at least $(2-\varepsilon)t$. Thus a universal constant
in Theorem~\ref{thm:main} must be at least two. Our constants are not
optimized.

The fractional coloring theorem of Reed and Seymour
\cite{ReedSeymour1998} gives $\chi_f(G)\le2h(G)$, where fractional
coloring permits weighted stable sets in place of color classes. It
provides a linear supply of stable sets but does not directly give a
linear ordinary or list coloring: naive repeated sampling introduces a
factor depending on the graph order. Delcourt and Postle
\cite{DelcourtPostle2021} obtained the ordinary bound
$O(t\log\log t)$ for graphs excluding a $K_t$ minor and reduced a linear
ordinary bound to coloring suitably small subgraphs. Liu and Luo
\cite{LiuLuo2026} subsequently proved the ordinary bound
$O(t\log\log\log t)$ and developed bounded stable-set sampling that we
use below.\footnote{Liu and Luo explicitly acknowledge assistance from
ChatGPT Codex; their declaration does not identify an underlying model
version.} These ordinary-coloring statements do not by themselves give
colorings from arbitrary vertex lists.

For list coloring, Postle proved an $O(t(\log\log t)^6)$ bound for
$K_t$-minor-free graphs \cite[Theorem~1.8]{Postle2020}. His separability
and woven-minor constructions provide the structural framework for our
passage from small graphs to arbitrary orders. Delcourt and Postle outlined
an $O(t(\log\log t)^2)$ list bound in the first version of their paper
\cite[Section~6]{DelcourtPostle2021v1}, leaving the details to the reader.
In subsequent work posted on October~1, 2026, Gu and Xu
\cite[Theorem~1.3]{GuXu2026} proved an $O(t\log\log t)$ list bound.
Connectivity extraction by the constraint-budget method of Gir\~ao and Narayanan
\cite{GiraoNarayanan2022}, and the highly linked graph theorem of
Bollob\'as and Thomason \cite{BollobasThomason1996}, are further key
antecedents. We reproduce the list transfer, density, path accounting,
and recursion arguments in the forms required here.

Two difficulties remain after these reductions. Even on a small graph,
a sampled stable set cannot be given a single color unless that color is
available at every vertex receiving it. In a graph of unrestricted order,
the paths used to assemble a minor may themselves contain arbitrarily
many vertices. Their deletion must therefore be charged through their
structure, rather than through their order. The proof addresses these
two difficulties separately.

\subsection{Proof strategy}

First, Proposition~\ref{prop:small-list} proves that a $K_t$-minor-free
graph on at most $t^{11/10}$ vertices is $Lt$-choosable, for an absolute
$L$ and all sufficiently large $t$. We partition the union of all lists
into bins, assigning each color to a single bin throughout the graph.
A vertex is eligible for a bin precisely when its own list contains a
color in that bin. A random partition gives every vertex enough eligible
bins at every stage, before any remaining graph is chosen.
Lemma~\ref{small:eligible-thinning} then repeatedly halves an upper bound
on the fractional chromatic number using sampled stable sets restricted to eligible
vertices. Vertices assigned to the same bin lie in one stable set, so
each can choose its own allowable color from that bin; colors in different
bins cannot conflict.

To prove the thinning lemma, suppose every sampling outcome leaves a
fractional chromatic number greater than half the current bound. A deterministic separator
decomposition finds a fixed highly connected subgraph in which a
remainder still has large expected fractional chromatic number. Independent
repetitions produce disjoint private vertex sets, each supporting a clique
model. A small, fixed family of matching tests shows that the complement
of the union of the sampled remainders remains highly connected.
Connecting equally labeled bags through this complement produces the
excluded clique minor. The proof uses bounded stable-set sampling and
ordinary-chromatic connectivity extraction, with the full argument given
in Section~\ref{small:section}. Summing the costs of successive halving
stages requires only $O(t)$ bins.

Second, Lemma~\ref{lem:list-extraction} extracts highly connected
subgraphs with a linear loss of list chromatic number. The additive
list-deletion estimate, Lemma~\ref{lem:additive-deletion}, reserves
colors for a small set and controls its interaction with all remaining
vertices. Its probabilistic estimate tests only polynomially many
exceptional vertices, independently of the ambient graph order and
palette size. Minimum-length path systems give the sparsity needed to
apply this estimate even when the paths are long. The density argument
also supplies small highly connected subgraphs in which clique models
can be built while connecting paths are rerouted.

The sequential construction in Lemma~\ref{lem:separability} then shows
that a graph of sufficiently large list chromatic number either contains
the desired clique minor or splits into two disjoint subgraphs retaining
nearly all of that list chromatic number. At each step, one small district
supplies new adjacencies between model bags and a separate junction joins
the paths carrying those bags forward. The deletion estimate applies to
the complete history of original minimizing paths, including vertices
no longer used by the growing model.

Finally, Theorem~\ref{thm:recursive-woven} builds a clique model whose
bags contain prescribed roots, together with paths for prescribed
terminal pairs. After reserving the paths for the terminal pairs, the
remaining task is a rooted model inside a fresh region. Separability
supplies three disjoint districts for recursive calls. Partitioning the
root labels into three blocks and assigning one district to each pair
of blocks reduces each call to about two thirds as many roots, while
every pair of labels still occurs together in one district. Connectivity
extraction and the case with no prescribed pairs complete the theorem.

Section~\ref{sec:preliminaries} fixes the external inputs and
probabilistic tools. Section~\ref{small:section} proves the small-graph
list bound, and Section~\ref{sec:list-transfer} establishes list
extraction and deletion. Section~\ref{sec:density} constructs small
highly connected districts; Section~\ref{sec:routing} develops their
routing interfaces. Sections~\ref{sec:separability} and
\ref{sec:recursion} prove separability and the final recursive theorem.
The only structural results used without proof are the fractional minor
bound, ordinary-chromatic connectivity extraction, and the terminal
grouping theorem, whose precise statements follow.
